\documentclass{article}
\usepackage[UTF8]{ctex}
\usepackage{amsmath,amssymb,amsfonts}
\usepackage{enumitem}
\usepackage{geometry}
\geometry{a4paper, margin=1in}
% === 错题本排版优化（不参与题目标识） ===
\setlength{\parindent}{0pt}
\setlength{\parskip}{0.35em plus 0.10em minus 0.08em}
\setlength{\emergencystretch}{3em}
\setlist{itemsep=0.25em, topsep=0.40em, parsep=0pt}
\allowdisplaybreaks[2]
% === 错题本排版优化结束 ===
\begin{document}
\section*{哈尔滨工业大学2025年数学分析真题}
\section*{82．哈尔滨工业大学2025年数学分析试卷}

\subsection{题目1}
1．判断题。正确的给出证明，错误的给出反例。

\begin{itemize}
  \item[（1）] 设 \(b_{n}=\frac{a_{1}+a_{2}+\cdots+a_{n}}{n}\) ，若数列 \(\left\{b_{n}\right\}\) 收敛，则 \(\left\{a_{n}\right\}\) 收敛。
  \item[（2）] 两个在 \(x_{0}\) 附近无界的函数之积仍为无界函数。
  \item[（3）] 若 \(y=f(x)\) 在点 \(x_{0}\) 不连续，\(u=g(y)\) 在点 \(y_{0}=f\left(x_{0}\right)\) 不连续，则复合函数 \(u=g(f(x))\) 在点\(x_{0}\) 不连续．
  \item[（4）] 一元函数的定积分，若 \(|f(x)|\) 可积，则 \(f(x)\) 可积．
  \item[（5）] 若二元函数 \(f(x, y)\) 在点 \(\left(x_{0}, y_{0}\right)\) 的两个二次极限都存在，则 \(f(x, y)\) 在 \(\left(x_{0}, y_{0}\right)\) 的二重极限也存在
\end{itemize}

\subsection{题目2}
2．解答如下问题：

\begin{itemize}
  \item[（1）] 若函数 \(f(x)\) 在 \((a, b)\) 上一致连续，证明：\(f(x)\) 在 \((a, b)\) 上有界．
  \item[（2）] 用两种方法证明函数 \(y=\ln x\) 在 \((0,1)\) 上不一致连续．
\end{itemize}

\subsection{题目3}
3．设函数 \(f(x, y)= \begin{cases}\left(x^{2}+y^{2}\right)^{\alpha} \sin \frac{1}{x^{2}+y^{2}}, & x^{2}+y^{2} \neq 0 ; \\ 0, & x^{2}+y^{2}=0 .\end{cases}\)

\begin{itemize}
  \item[（1）] 求 \(\alpha\) 的范围，使函数 \(f(x, y)\) 在 \((0,0)\) 连续．
  \item[（2）] 求 \(\alpha\) 的范围，使函数 \(f(x, y)\) 在 \((0,0)\) 可微．
  \item[（3）] 求 \(\alpha\) 的范围，使函数 \(f(x, y)\) 在 \((0,0)\) 的偏导数连续．
\end{itemize}

\subsection{题目4}
4．分别从微积分的角度，用五种方法证明：\(e^{x} \geq 1+x\) ．

\subsection{题目5}
5．若对 \(\forall x \in[a, b]\) ，函数 \(f(x)\) 满足 \(f(x) \geq 0, f^{\prime \prime}(x) \leq 0\) ．求证：\(f(x) \leq \frac{2}{b-a} \int_{a}^{b} f(x) \mathrm{d} x\) ．

\subsection{题目6}
6．设函数 \(f(x)\) 在 \((a, b)\) 上连续且无极大值点，\(\left\{x_{n}\right\} \subset(a, b)\) 满足 \(\lim _{n \rightarrow \infty} f\left(x_{n}\right)=\inf _{a<x<b} f(x)\) ．

\begin{itemize}
  \item[（1）] 用肯定语言叙述 \(\lim _{n \rightarrow \infty} x_{n} \neq \bar{x}\) 的 \(\varepsilon-N\) 定义。
  \item[（2）] 若 \(\lim _{n \rightarrow \infty} x_{n}=x_{0}\) ，是否有 \(x_{0} \in(a, b)\) ？如果有，请加以证明，否则请举反例．
  \item[（3）] 证明 \(\lim _{n \rightarrow \infty} x_{n}\) 存在．
\end{itemize}

\subsection{题目7}
7．设正项级数 \(\sum_{n=1}^{\infty} \frac{1}{p_{n}}\) 收敛．记 \(A_{n}=\sum_{k=1}^{n} p_{k}, B_{n}=\sum_{k=1}^{n} \frac{k^{2}}{A_{k}^{2}}, C_{n}=\sum_{k=1}^{n} \frac{1}{p_{k}}, S_{n}=\sum_{k=1}^{n} \frac{k}{A_{k}}\) ．证明：

\begin{itemize}
  \item[（1）] \(B_{n}<\frac{5}{p_{1}}+2 S_{n}+C_{n}\) ．
  \item[（2）] \(S_{n}<\sqrt{B_{n} C_{n}}\) ．
  \item[（3）] \(\sum_{n=1}^{\infty} \frac{n}{p_{1}+p_{2}+\cdots+p_{n}}\) 收敛．
\end{itemize}

\subsection{题目8}
8．证明：

\begin{itemize}
  \item[（1）] \(\sum_{n=1}^{\infty} x^{n} \ln x\) 在 \((0,1]\) 上不一致收敛。
  \item[（2）] \(\sum_{n=1}^{\infty} x^{n} \frac{\ln x}{1+|\ln (-\ln x)|}\) 在 \((0,1)\) 上一致收敛．
\end{itemize}

\subsection{题目9}
9．设 \(I(k, a)=\int_{0}^{+\infty} e^{-k x} \frac{\sin a x}{x} \mathrm{~d} x\) ．\\
（1）证明：对固定的 \(k \in[0,+\infty)\) ，积分 \(I(k, a)\) 关于 \(a\) 在 \(|a| \geq \delta(\delta>0)\) 上一致收敛．对固定的 \(a \neq 0\) ，积分 \(I(k, a)\) 关于 \(k\) 在 \([0,+\infty)\) 上一致收敛。\\
（2）计算积分 \(I=\int_{0}^{+\infty} \frac{\sin x}{x} \mathrm{~d} x\) 的值．

\subsection{题目10}
10．计算：\\
（1）\(\oint_{C} \frac{y}{x^{2}+4 y^{2}} \mathrm{~d} x-\frac{x}{x^{2}+4 y^{2}} \mathrm{~d} y\) ，其中 \(C\) 是圆周 \(x^{2}+y^{2}=1\) ，取逆时针方向．\\
（2） \(\iint_{S}(y-z) \mathrm{d} y \mathrm{~d} z+(z-x) \mathrm{d} z \mathrm{~d} x+(x-y) \mathrm{d} x \mathrm{~d} y\) ，其中 \(S\) 是圆锥 \(x^{2}+y^{2}=z^{2}(0 \leq z \leq 2)\) 的外侧。\\
（3）\(\oint_{C} y^{3} \mathrm{~d} x+z^{3} \mathrm{~d} y+x^{3} \mathrm{~d} z\) ，其中 \(C\) 为球面 \(x^{2}+y^{2}+z^{2}=4\) 与平面 \(x+y+z=0\) 的交线，从 \(z\) 轴正向看为逆时针方向。
\end{document}
